\documentclass[11pt]{article}

\usepackage[a4paper,margin=1in]{geometry}
\usepackage{amsmath,amssymb,mathtools}
\usepackage{enumitem}
\usepackage{booktabs}
\usepackage{array}
\usepackage{xcolor}
\usepackage{listings}
\usepackage[hidelinks]{hyperref}

\lstdefinestyle{plaincode}{
  basicstyle=\ttfamily\small,
  columns=fullflexible,
  keepspaces=true,
  breaklines=true,
  frame=single,
  rulecolor=\color{black!20}
}
\lstset{style=plaincode}

\newcommand{\QQ}{\mathbb{Q}}
\newcommand{\ZZ}{\mathbb{Z}}
\newcommand{\FF}{\mathbb{F}}
\newcommand{\ii}{\mathrm{i}}
\newcommand{\supp}{\operatorname{supp}}
\newcommand{\rank}{\operatorname{rank}}
\newcommand{\xor}{\mathbin{\mathtt{xor}}}

\title{Documentation for a Sparse Algebra of Gaussian Rational Radical Numbers with \texttt{qi\_algebra.py}}
\author{}
\date{}

\begin{document}
\maketitle

\section{Purpose and scope}

This document describes the mathematical model and implementation strategy for the \texttt{QINumber} class whose elements are finite sums of the form
\begin{equation}
  A = \sum_{j = 1}^{N_{\rm terms}} \frac{p_j + \ii q_j}{d_j}\sqrt{r_j},
\end{equation}
where
\[
  p_j,q_j,d_j,r_j \in \ZZ, \qquad d_j>0,
\]
and where the sum is finite. The coefficient
\[
  \frac{p_j+\ii q_j}{d_j}
\]
is a Gaussian rational number, i.e. an element of \(\QQ(\ii)\). The radical part \(\sqrt{r_j}\) is normalized to be a product of prime-factors i.e.
\begin{equation}
\label{eq:1}
r_j = 2^{n_{0}} \times 3^{n_1} \times \dots \times \mathfrac{P}_j^{n_{j-1}} \times \dots,
\end{equation}
where $P_j$ is the $j$-th prime and $n_{j} \in \{0, 1\}$ for all $j$. For ease of implementation the radicands are indexed by integers see Sec.~\ref{squareclass_encoding}.

The class implements exact arithmetic for addition, subtraction, negation, complex conjugation, multiplication, inversion, division, integer powers, equality, and normalization.

The practical design goal is:
\begin{quote}
Represent each number as a sparse dictionary from a radical squareclass index to a normalized Gaussian rational coefficient.
\end{quote}

The mathematical design goal is:
\begin{quote}
Work inside finite extensions of \(\QQ(\ii)\) generated by square roots: \(\QQ(\ii, \sqrt{r_1}, \dots )\) while using only integer arithmetic, sparse dictionaries, and binary vector operations over \(\FF_2\).
\end{quote}

\section{Canonical representation}

\subsection{High-level representation}

An element is stored as a dictionary
\begin{equation}
  \texttt{data}: \texttt{radicand\_index} \longmapsto [p,q,d].
\end{equation}
The entry
\[
  i \mapsto [p,q,d]
\]
represents the term
\begin{equation}
  \frac{p+\ii q}{d}\sqrt{R_i},
\end{equation}
where \(R_i\) is the squarefree radicand encoded by the integer index \(i\).

The zero number is represented by the empty dictionary:
\begin{equation}
  0 \equiv \{\}.
\end{equation}

A pure Gaussian rational
\[
  \frac{p+\ii q}{d}
\]
is represented by
\begin{equation}
  \{0:[p,q,d]\},
\end{equation}
provided \((p,q)\neq (0,0)\). The index \(0\) corresponds to the empty squareclass, i.e. \(\sqrt{1}=1\).

\subsection{Encoding squareclasses by bitmasks}
\label{squareclass_encoding}

Every integer radicand is reduced to its squarefree part. If
\begin{equation}
  r = \pm \prod_{k\geq 0} p_k^{n_k},
\end{equation}
where \(p_0=2,p_1=3,p_2=5,\ldots\) are the primes, then only the parities
\begin{equation}
  \epsilon_k = n_k \bmod 2
\end{equation}
matter for the squareclass.

The squareclass is encoded by the integer bitmask
\begin{equation}
  i = \sum_{k\geq 0} \epsilon_k 2^k.
\end{equation}
Thus:
\begin{align*}
  \sqrt{1} &\leftrightarrow 0,\\
  \sqrt{2} &\leftrightarrow 1,\\
  \sqrt{3} &\leftrightarrow 2,\\
  \sqrt{5} &\leftrightarrow 4,\\
  \sqrt{6}=\sqrt{2\cdot 3} &\leftrightarrow 1+2=3,\\
  \sqrt{30}=\sqrt{2\cdot3\cdot5} &\leftrightarrow 1+2+4=7.
\end{align*}

The bitmask should be understood as a vector in \(\FF_2^m\), where each bit records whether a prime appears with odd exponent.

\subsection{Multiplication of radical indices}

Multiplication of squareclasses is bitwise XOR:
\begin{equation}
  \sqrt{R_a}\sqrt{R_b}
  = \gamma(a,b)\sqrt{R_{a\xor b}},
\end{equation}
where \(a\xor b\) is bitwise XOR. The scalar \(\gamma(a,b)\in \ZZ\) accounts for square factors produced by the product.

If the stored indices already correspond to squarefree radicands \(R_a,R_b\), then
\begin{equation}
  R_aR_b = G(a,b)^2 R_{a\xor b},
\end{equation}
so
\begin{equation}
  \sqrt{R_a}\sqrt{R_b}=G(a,b)\sqrt{R_{a\xor b}}.
\end{equation}
Here \(G(a,b)\) is the product of primes whose bits occur in both \(a\) and \(b\). Equivalently,
\begin{equation}
  G(a,b)=R_{a\,\&\,b},
\end{equation}
where \(a\,\&\,b\) is bitwise AND, interpreted as the product of the corresponding primes.

Example:
\begin{equation}
  \sqrt{6}\sqrt{10}
  = \sqrt{60}
  = 2\sqrt{15}.
\end{equation}
The indices are
\[
  6\leftrightarrow 3=(011)_2, \qquad 10\leftrightarrow 5=(101)_2.
\]
Then
\[
  3\xor 5 = 6=(110)_2,
\]
corresponding to \(15=3\cdot5\), and
\[
  3\,\&\,5 = 1,
\]
corresponding to the shared prime \(2\), so the scalar factor is \(2\).

\section{Normalization invariants}

After every public operation, the object should be normalized. The normal form has the following invariants.

\begin{enumerate}[label=\textbf{N\arabic*.}]
  \item The zero element is exactly the empty dictionary \(\{\}\).
  \item No stored coefficient is zero. That is, there is no entry \([0,0,d]\).
  \item Every denominator is positive: \(d>0\).
  \item Each coefficient \([p,q,d]\) is reduced by
  \begin{equation}
    g = \gcd(p,q,d),
  \end{equation}
  replacing
  \begin{equation}
    [p,q,d] \mapsto [p/g,q/g,d/g].
  \end{equation}
  \item Like radical indices are combined. There is at most one dictionary entry for each squareclass index.
  \item All radicand indices are squarefree bitmask indices. No prime appears more than once in a stored radical basis element.
\end{enumerate}

Equality is then dictionary equality:
\begin{equation}
  A=B \quad\Longleftrightarrow\quad \operatorname{normalize}(A).\texttt{data}=\operatorname{normalize}(B).\texttt{data}.
\end{equation}

\section{Gaussian rational coefficient arithmetic}

The coefficient associated to a radical index is stored as
\begin{equation}
  [p,q,d] \quad\leftrightarrow\quad \frac{p+\ii q}{d}.
\end{equation}

The basic coefficient operations are:
\begin{align*}
  \frac{p_1+\ii q_1}{d_1} + \frac{p_2+\ii q_2}{d_2}
  &= \frac{(p_1d_2+p_2d_1)+\ii(q_1d_2+q_2d_1)}{d_1d_2},\\
  -\frac{p+\ii q}{d}
  &=\frac{-p-\ii q}{d},\\
  \overline{\frac{p+\ii q}{d}}
  &=\frac{p-\ii q}{d},\\
  \frac{p_1+\ii q_1}{d_1}\frac{p_2+\ii q_2}{d_2}
  &=\frac{(p_1p_2-q_1q_2)+\ii(p_1q_2+q_1p_2)}{d_1d_2}.
\end{align*}

A nonzero Gaussian rational coefficient is inverted by
\begin{equation}
  \left(\frac{p+\ii q}{d}\right)^{-1}
  = \frac{d(p-\ii q)}{p^2+q^2}.
\end{equation}

\section{Addition, subtraction, negation, and conjugation}

\subsection{Addition}

Addition is dictionary merging.

Given two normalized objects \(A\) and \(B\):
\begin{enumerate}[label=\alph*)]
  \item Copy one dictionary into the result.
  \item For each term of the other dictionary:
  \begin{enumerate}[label=\roman*)]
    \item if the radical index is absent, insert the coefficient;
    \item if the radical index is present, add the Gaussian rational coefficients.
  \end{enumerate}
  \item Normalize the result, removing zero coefficients.
\end{enumerate}

\subsection{Subtraction}

Subtraction is implemented as
\begin{equation}
  A-B = A+(-B).
\end{equation}
Negation sends every coefficient
\begin{equation}
  [p,q,d]\mapsto [-p,-q,d].
\end{equation}

\subsection{Complex conjugation}

Complex conjugation acts only on \(\ii\), not on the formal real radicals:
\begin{equation}
  \overline{\sqrt{R_i}}=\sqrt{R_i}.
\end{equation}
Therefore each coefficient changes by
\begin{equation}
  [p,q,d]\mapsto [p,-q,d].
\end{equation}

\section{Multiplication}

Multiplication is term-by-term distribution followed by collection.

If
\begin{equation}
  A=\sum_a c_a\sqrt{R_a},\qquad
  B=\sum_b e_b\sqrt{R_b},
\end{equation}
then
\begin{equation}
  AB=\sum_{a,b} c_a e_b G(a,b)\sqrt{R_{a\xor b}}.
\end{equation}

The implementation loops over all pairs of dictionary entries:
\begin{lstlisting}
result = {}
for a, ca in A.data.items():
    for b, cb in B.data.items():
        g = shared_square_factor(a, b)      # product of common primes
        k = a xor b                         # new squareclass index
        coeff = ca * cb * g
        add coeff to result[k]
normalize(result)
\end{lstlisting}

For sparse objects with \(M\) and \(N\) terms, raw multiplication takes \(O(MN)\) coefficient products before collection.

\section{Inversion}

\subsection{Why inversion is always possible for nonzero elements}

Let
\begin{equation}
  A=\sum_j c_j\sqrt{R_j}, \qquad c_j\in\QQ(\ii).
\end{equation}
The finitely many radicals appearing in \(A\) generate a finite field extension
\begin{equation}
  K=\QQ(\ii,\sqrt{R_1},\ldots,\sqrt{R_n}).
\end{equation}
If \(A\neq 0\), then \(A^{-1}\in K\). Therefore the inverse can again be represented as a finite sum of Gaussian rational coefficients times square-root basis elements.

The implementation does not need general field machinery. It computes the inverse by repeated rationalization of the denominator.

\subsection{Single-term inverse}

If
\begin{equation}
  A=c\sqrt{R_i},\qquad c\in\QQ(\ii),\quad c\neq0,
\end{equation}
then
\begin{equation}
  A^{-1}=\frac{1}{cR_i}\sqrt{R_i}.
\end{equation}
This follows from
\begin{equation}
  \frac{1}{c\sqrt{R_i}}=\frac{\sqrt{R_i}}{cR_i}.
\end{equation}
If \(i=0\), this is just inversion of a Gaussian rational.

\subsection{Two-term inverse}

If
\begin{equation}
  A=x+y,
\end{equation}
where \(x\) and \(y\) are two radical terms, then multiply numerator and denominator by
\begin{equation}
  x-y.
\end{equation}
Then
\begin{equation}
  A^{-1}=\frac{x-y}{x^2-y^2}.
\end{equation}
The denominator has fewer radical components, and in the two-term case it becomes a Gaussian rational after collection.

\subsection{General inverse by repeated denominator rationalization}

For more than two terms, the inverse is computed by maintaining
\begin{equation}
  \frac{N}{D}=\frac{1}{A},
\end{equation}
initialized as
\begin{equation}
  N=1,\qquad D=A.
\end{equation}
At every step, choose a squareclass direction and split
\begin{equation}
  D=E+\theta O,
\end{equation}
where \(\theta\) is a chosen radical generator and neither \(E\) nor \(O\) contains \(\theta\). Then multiply both numerator and denominator by
\begin{equation}
  E-\theta O.
\end{equation}
The new denominator is
\begin{equation}
  D'=(E+\theta O)(E-\theta O)=E^2-\theta^2O^2.
\end{equation}
The chosen squareclass direction has been eliminated from the denominator.

Repeat until the denominator is a Gaussian rational. Then invert that final Gaussian rational and multiply it into the numerator.

\section{The square-space over \texorpdfstring{\(\FF_2\)}{F2}}

\subsection{Support set}

The radical indices currently appearing in the denominator form a finite support set
\begin{equation}
  S = \supp(D)\subseteq \FF_2^m.
\end{equation}
Each index is a bit vector. The rank of the denominator is
\begin{equation}
  \rho=\rank_{\FF_2}\langle S\rangle.
\end{equation}
This rank is the number of independent square-root directions that still need to be eliminated.

The minimum possible number of rationalization steps is \(\rho\), assuming each step eliminates one independent direction.

\subsection{Cuts are linear functionals}

A rationalization step is determined by a nonzero linear functional
\begin{equation}
  \lambda:\FF_2^m\to\FF_2.
\end{equation}
It splits the support into two sets:
\begin{align}
  S_0&=\{v\in S:\lambda(v)=0\},\\
  S_1&=\{v\in S:\lambda(v)=1\}.
\end{align}
This corresponds to writing
\begin{equation}
  D=E+\theta O.
\end{equation}

The conjugating factor is
\begin{equation}
  E-\theta O.
\end{equation}

\subsection{Support update after a cut}

The new denominator support is contained in
\begin{equation}
  S' = (S_0+S_0)\cup(S_1+S_1),
\end{equation}
where addition is XOR in \(\FF_2^m\). Equivalently,
\begin{equation}
  S' = \{u\xor v:u,v\in S_0\}\cup\{u\xor v:u,v\in S_1\}.
\end{equation}
Actual coefficient arithmetic may cause additional cancellations, so this set is the structural support before accidental zero coefficients are removed.

The exact support can be computed by pairwise XORs inside each side of the cut, while simultaneously computing the coefficients of \(E^2-\theta^2O^2\).

\subsection{Cost of a candidate cut}

Let
\begin{equation}
  e=|S_0|,\qquad o=|S_1|,\qquad e+o=|S|=N.
\end{equation}
Computing the denominator by symmetric squaring takes approximately
\begin{equation}
  C(\lambda)=\frac{e(e+1)}{2}+\frac{o(o+1)}{2}
\end{equation}
term multiplications.

For fixed \(N\), this expression is minimized when \(e\) and \(o\) are as close as possible. Therefore the first heuristic is:
\begin{quote}
Prefer cuts whose support is close to half of the current terms.
\end{quote}

A stronger one-step heuristic is:
\begin{quote}
Among balanced cuts, prefer the cut with the smallest resulting support
\[
  |(S_0+S_0)\cup(S_1+S_1)|.
\]
\end{quote}

\section{Greedy inversion strategy}

\subsection{Precomputation}

At the beginning of inversion:
\begin{enumerate}[label=\arabic*.]
  \item Extract the denominator support \(S\).
  \item Compute the \(\FF_2\)-rank \(\rho\).
  \item Compute a basis for the active square-space.
  \item Compute a corresponding dual basis of linear functionals, or otherwise generate candidate functionals \(\lambda\).
  \item For each candidate, estimate the split size and optionally the exact next support size.
\end{enumerate}

In implementation language, the ``orthogonal basis'' is best interpreted as a dual system of parity tests. A move is not chosen by a radical term directly; it is chosen by a parity test on the radical bitmasks.

\subsection{Candidate ordering}

A cheap candidate starts from prime-factor columns. For a prime column \(k\), define
\begin{equation}
  B_k=\{v\in S:v_k=1\}.
\end{equation}
Its support size is \(|B_k|\), and the induced split is
\begin{equation}
  |B_k| \quad\text{versus}\quad N-|B_k|.
\end{equation}
A good single-column candidate has \(|B_k|\) close to \(N/2\).

More general candidates are XOR combinations of prime columns. If the current candidate support is \(B\) and a new prime column is \(C\), then adding that prime to the parity test replaces \(B\) by the symmetric difference
\begin{equation}
  B\triangle C.
\end{equation}
The size can be updated by
\begin{equation}
  |B\triangle C|=|B|+|C|-2|B\cap C|.
\end{equation}
This makes it cheap to greedily improve a candidate by adding columns that move the support closer to half of \(N\).

A practical ordering is:
\begin{enumerate}[label=\arabic*.]
  \item Rank prime columns by closeness of support to \(N/2\).
  \item Use the best columns as seeds.
  \item If a seed has support larger than \(N/2\), try adding other columns by XOR to bring the support closer to \(N/2\).
  \item Also allow adding columns to smaller seeds if it improves balance.
  \item For the resulting candidates, compute the actual denominator cost and, if desired, exact next support.
\end{enumerate}

This is a heuristic. It is not a proof of global optimality. It is intended to minimize immediate multiplication cost and reduce expression swell.

\subsection{Dynamic programming versus greedy search}

The exact minimal number of rationalization steps is the rank \(\rho\). The exact minimum total multiplication cost over all possible orders is a shortest-path or dynamic-programming problem over denominator supports.

A denominator-only dynamic program can be written as:
\begin{lstlisting}
best(S):
    if S == {0}:
        return 0
    if S in memo:
        return memo[S]

    ans = infinity
    for lambda in candidate_functionals(span(S)):
        S0 = {v in S : dot(lambda, v) == 0}
        S1 = {v in S : dot(lambda, v) == 1}

        e = len(S0)
        o = len(S1)
        step_cost = e*(e+1)//2 + o*(o+1)//2

        Snext = xor_pair_support(S0) union xor_pair_support(S1)
        ans = min(ans, step_cost + best(Snext))

    memo[S] = ans
    return ans
\end{lstlisting}

In practice, the search complexity depends more on \(\rho\) than on the number of terms \(N\). There are up to
\begin{equation}
  2^\rho-1
\end{equation}
nonzero linear functionals at rank \(\rho\). Therefore the practical policy is:

\begin{center}
\begin{tabular}{>{\raggedright\arraybackslash}p{0.25\linewidth} >{\raggedright\arraybackslash}p{0.65\linewidth}}
\toprule
Regime & Suggested strategy \\
\midrule
\(N\leq 3\) & Direct formulas or hand-simplified rationalization. \\
\(3<N\leq 10\) & Greedy is usually cheaper than setting up full dynamic programming. \\
\(10<N\leq 50\), small \(\rho\) & Dynamic programming or bounded lookahead can be worthwhile. \\
\(10<N\leq 50\), large \(\rho\) & Greedy with exact one-step support calculation, or beam search. \\
\(N>50\) & Greedy heuristics; possibly beam search if expression swell is severe. \\
\bottomrule
\end{tabular}
\end{center}

\section{Special handling for three terms}

For three-term denominators, the implementation uses the general rationalization idea but simplifies the cut selection by hand.

The goal is to choose a split that does not separate terms whose radicands share a common factor in a way that immediately recreates unwanted radicals through cross-products. In bitmask language, the cut should avoid an obviously bad split when two radicals have a large overlap in their prime-factor parity vectors.

This three-term rule is a small-case heuristic. It avoids the overhead of building the full square-space search machinery for a denominator so small that the general optimization machinery is unlikely to pay for itself.

\section{Division}

Division is implemented as multiplication by the inverse:
\begin{equation}
  \frac{A}{B}=A B^{-1}, \qquad B\neq0.
\end{equation}
The divisor is inverted using the inversion logic above, and the result is multiplied by the dividend.

\section{Integer powers}

The method computes \(A^n\) for integer \(n\).

\subsection{Negative, zero, and small positive exponents}

\begin{enumerate}[label=\arabic*.]
  \item If \(n<0\), first compute \(A^{-1}\), then compute \((A^{-1})^{-n}\).
  \item If \(n=0\), return \(1\), except in the undefined case \(0^0\) where \texttt{DomainError} is raised.
  \item If \(n=1\), return \(A\).
  \item If \(n=2\) or \(n=3\), use direct multiplication. The exponent is too small for optimization overhead to be useful.
\end{enumerate}

\subsection{Single-term powers}

If
\begin{equation}
  A=c\sqrt{R_i},
\end{equation}
then even powers become Gaussian rationals times powers of \(R_i\), and odd powers are that even-power result times \(A\). Explicitly:
\begin{align*}
  A^{2k} &= c^{2k} R_i^k,\\
  A^{2k+1} &= c^{2k+1} R_i^k\sqrt{R_i}.
\end{align*}
Because the result has at most one radical term, this case can be exponentiated directly using coefficient arithmetic and integer powers.

\subsection{Rank-one square-space powers}

If the support of \(A\) has rank one over \(\FF_2\), then \(A\) lies in a quadratic extension of \(\QQ(\ii)\). Therefore \(A\) satisfies a quadratic equation
\begin{equation}
  a x^2 + b x + c = 0,
\end{equation}
with coefficients in \(\QQ(\ii)\). Rearranging gives
\begin{equation}
  x^2 = -\frac{b}{a}x - \frac{c}{a}.
\end{equation}
Thus every power of \(x=A\) reduces to a linear expression
\begin{equation}
  x^n = \alpha_n x + \beta_n.
\end{equation}

The coefficients \(\alpha_n,\beta_n\in\QQ(\ii)\) are computed by matrix exponentiation of the recurrence induced by
\begin{equation}
  x^2 = ux+v,
\end{equation}
where
\begin{equation}
  u=-\frac{b}{a}, \qquad v=-\frac{c}{a}.
\end{equation}

Multiplication by \(x\) acts on the vector representation \(\alpha x+\beta\) by
\begin{equation}
  x(\alpha x+\beta)=\alpha x^2+\beta x
  =\alpha(ux+v)+\beta x
  =(\alpha u+\beta)x+\alpha v.
\end{equation}
Therefore the transition is
\begin{equation}
  \begin{pmatrix}
  \alpha'\\
  \beta'
  \end{pmatrix}
  =
  \begin{pmatrix}
  u & 1\\
  v & 0
  \end{pmatrix}
  \begin{pmatrix}
  \alpha\\
  \beta
  \end{pmatrix}.
\end{equation}
Starting from
\begin{equation}
  x^1 = 1\cdot x + 0,
\end{equation}
one obtains \(x^n\) by exponentiating this \(2\times2\) matrix.

\subsection{Rank two and higher powers}

Rank-two expressions could in principle be optimized by using the degree-four algebra generated by two independent square-root directions. However, the implementation does not currently do this because:
\begin{enumerate}[label=\arabic*.]
  \item the implementation complexity is higher;
  \item the practical computational benefit is unclear;
  \item multiplication is already exact and simple for moderate exponents.
\end{enumerate}

For general positive exponents outside the optimized cases, the implementation uses repeated multiplication.

\section{Equality}

Equality is defined by normalized dictionary equality.

Because all terms are collected by squareclass index and all coefficients are reduced, two mathematically equal numbers have the same normalized dictionary representation, assuming all construction paths respect the same squarefree-index convention.

Thus:
\begin{equation}
  A==B
\end{equation}
means:
\begin{equation}
  \operatorname{normalize}(A).\texttt{data}
  ==
  \operatorname{normalize}(B).\texttt{data}.
\end{equation}

\section{Implementation-oriented summary}

The object represents elements of a sparse \(\QQ(\ii)\)-vector space with basis elements \(\sqrt{R_i}\), where \(i\) is a bitmask encoding the squarefree prime support of \(R_i\).

\begin{center}
\begin{tabular}{>{\raggedright\arraybackslash}p{0.32\linewidth} >{\raggedright\arraybackslash}p{0.58\linewidth}}
\toprule
Concept & Implementation meaning \\
\midrule
Gaussian rational coefficient & Triple \([p,q,d]\) representing \((p+\ii q)/d\). \\
Radical basis element & Integer bitmask representing squarefree prime support. \\
Addition & Merge dictionaries and add coefficients with equal keys. \\
Negation & Negate \(p\) and \(q\) in every coefficient. \\
Complex conjugation & Negate only \(q\). \\
Multiplication & Pairwise term multiplication; new key is XOR of keys; coefficient gets square-factor correction. \\
Normalization & Reduce coefficient gcds, force positive denominators, combine keys, remove zero terms. \\
Inversion & Repeatedly rationalize denominator using \(D=E+\theta O\), multiply by \(E-\theta O\). \\
Rank & Number of independent square-root directions in the current support. \\
Cut & Linear parity test \(\lambda(v)\) on bitmask vectors. \\
Greedy cut score & Prefer balanced splits and smaller next support. \\
Division & Multiply by inverse. \\
Power & Special cases for negative, zero, small, single-term, and rank-one; otherwise repeated multiplication. \\
Equality & Normalized dictionary equality. \\
\bottomrule
\end{tabular}
\end{center}

\section{LLM-facing interpretation guide}

The code is not trying to implement general symbolic algebra. It implements one specific exact algebra:
\begin{equation}
  \QQ(\ii)\text{-linear combinations of square roots of squarefree integers.}
\end{equation}

The important facts for understanding the code are:

\begin{enumerate}[label=\arabic*.]
  \item The dictionary key is not the original radicand. It is the squarefree squareclass encoded as a prime-parity bitmask.
  \item Multiplying radicals means XORing keys, plus multiplying the coefficient by the product of shared primes.
  \item Coefficients are exact Gaussian rationals stored as integer triples.
  \item Normalization is essential. Without normalization, equality by dictionary comparison is invalid.
  \item Inversion works because finitely many square roots generate a finite field extension, and every nonzero field element has an inverse in that same field.
  \item The inversion algorithm is repeated rationalization, not approximate numerical inversion.
  \item The GF(2) rank measures how many independent radical directions remain in a denominator.
  \item A rationalization step is a parity cut of the current radical support.
  \item Greedy inversion is an optimization heuristic for reducing multiplication cost and expression swell. It is not guaranteed globally optimal.
  \item Rank-one powers are optimized because the element satisfies a quadratic equation, so all powers reduce to \(\alpha A+\beta\).
\end{enumerate}

These points should be enough to understand why the algorithms are correct and why the implementation uses bitmasks, XOR, dictionary collection, and repeated conjugation.

\section{Potential edge cases}

The implementation should explicitly handle or reject the following cases:

\begin{enumerate}[label=\arabic*.]
  \item Inversion of zero is undefined $\to$ \texttt{MathDomainError}.
  \item Division by zero is undefined $\to$ \texttt{MathDomainError}.
  \item \(0^0\) is undefined $\to$ \texttt{MathDomainError}.
  \item Denominators in coefficient triples must never be zero. Zero is represented as $0 = \{\}$ and zero terms from a non-zero \texttt{QINumber} are removed from its dictionary.
  \item Denominators should be normalized to positive values.
  \item Zero coefficients must be removed from the dictionary.
  \item Original radicands must be reduced to squarefree squareclass indices before storage.
  \item Negative radicands are not allowed. Attempts to pass negative radicands $\to$ \texttt{MathDomainError} with the justification that the sign of negative radicands is arbitrary, as  $\sqrt{-1} = \pm i$.

\end{enumerate}

\end{document}
